\subsection{\normalsize{Enfoques de compresión sin pérdida para datos de punto flotante}}

En simulación científica es frecuente trabajar con arreglos extensos de números de
punto flotante, cuyo volumen puede convertir el almacenamiento y el movimiento de
datos en una restricción de desempeño. En este contexto, la compresión sin pérdida
busca reducir el tamaño de la representación binaria sin alterar ningún bit de la
información original, de modo que la reconstrucción coincida exactamente con los datos
de entrada. Esta exigencia es especialmente relevante en escenarios donde los valores
comprimidos participan posteriormente en nuevas etapas de cálculo o sirven como
referencia numérica exacta.

A diferencia de la compresión de texto o de datos discretos, la compresión de punto
flotante presenta dificultades particulares. Muchos arreglos científicos contienen
redundancias asociadas a continuidad espacial, dependencia temporal o regularidades en
la representación binaria de signo, exponente y mantisa. Sin embargo, cuando esas
regularidades son débiles o no siguen una estructura local simple, la compresión se
vuelve más desafiante. Por esta razón, la literatura ha propuesto distintos enfoques que
no actúan todos del mismo modo: algunos se basan en predicción, otros en
transformaciones reversibles de la representación binaria y otros en reorganización
tipada de los datos antes de aplicar un códec general.

\subsubsection{FPC (Burtscher \& Ratanaworabhan)}

\textit{FPC} (\emph{Floating Point Compressor})\parencite{burtscher2009fpc} es un algoritmo
de compresión sin pérdida diseñado para secuencias lineales de valores de punto
flotante, especialmente en doble precisión. Su principio central consiste en predecir el
siguiente valor a partir del historial reciente y comprimir únicamente la diferencia entre el
valor real y el valor estimado.

Para ello, FPC emplea dos predictores derivados de esquemas tipo FCM/DFCM, cuyos
resultados se comparan con el dato real mediante una operación \texttt{XOR}. Cuando la
predicción es cercana, la diferencia resultante presenta una cantidad considerable de
bytes iniciales nulos, que pueden codificarse de forma compacta. El algoritmo almacena
entonces un pequeño código de control junto con los bytes residuales no nulos. De esta
manera, la compresión depende directamente de la capacidad del predictor para anticipar
la estructura local de la secuencia.

Este enfoque resulta especialmente adecuado cuando existe correlación entre valores
adyacentes o cuando la evolución del arreglo mantiene cierta regularidad. En cambio, su
efectividad disminuye cuando los datos no exhiben continuidad local suficiente para que
la predicción produzca residuos pequeños. Aun así, FPC constituye una referencia
importante en compresión sin pérdida de punto flotante porque muestra con claridad la
utilidad y las limitaciones de los esquemas predictivos sobre representaciones IEEE 754.

\subsubsection{FPZIP (Lindstrom \& Isenburg)}

\textit{FPZIP}\parencite{lindstrom2006fpzip} es un compresor sin pérdida orientado a arreglos
de punto flotante en contextos científicos. Su diseño parte de la idea de que muchos
conjuntos de datos numéricos presentan continuidad o dependencia local, de modo que
pueden describirse con mayor eficiencia a partir de errores de predicción que a partir de
los valores originales.

A diferencia de enfoques lineales simples, FPZIP incorpora esquemas de predicción
adaptados a la estructura de los datos, lo que le permite aplicarse a arreglos de
distintas dimensionalidades. Tras generar una estimación para cada valor, el algoritmo
codifica el residuo de predicción mediante un modelo entrópico diseñado para preservar
exactitud bit a bit sobre la representación en punto flotante. En este sentido, FPZIP no
requiere cuantización ni aproximación intermedia: la compresión sigue siendo
estrictamente reversible.

La relevancia de FPZIP radica en que explota de manera más
general la estructura local de arreglos científicos y combina predicción con codificación
entrópica especializada. Esto lo convierte en un referente importante para datos en los
que la redundancia no se manifiesta únicamente entre valores consecutivos, sino también
en relaciones espaciales o estructurales más amplias.

\subsubsection{Bitshuffle con LZ4 (Masui et al.)}

Una estrategia diferente consiste en no predecir los valores, sino transformar su
disposición binaria para hacer visibles regularidades que no son evidentes en el orden
original. Bajo esta idea, \textit{Bitshuffle}\parencite{masui2015bitshuffle} actúa como un filtro
reversible que reorganiza los bits de una secuencia de valores de punto flotante antes de
aplicar un códec de compresión general, usualmente LZ4.

El procedimiento puede entenderse como una transposición bit a bit: si se consideran
$N$ valores de $m$ bits, sus representaciones binarias forman una matriz de
$N \times m$, y Bitshuffle reordena los datos agrupando los bits según su posición. Esta
transformación puede revelar repeticiones o patrones en determinados planos de bits,
por ejemplo en exponentes o signos, que de otra forma quedarían dispersos entre los
bytes de cada valor. Una vez realizada la transposición, un códec LZ puede aprovechar
con mayor eficacia esas secuencias más homogéneas.

Este enfoque no depende de predicciones locales entre valores adyacentes, sino de la
posibilidad de exponer regularidades internas de la representación binaria. Por ello, su
interés en compresión científica es doble: por un lado, preserva completamente la
exactitud; por otro, ofrece una vía distinta para explotar redundancia en datos donde los
modelos predictivos no siempre resultan adecuados. En la práctica, Bitshuffle suele
entenderse mejor como una transformación previa al códec que como un compresor
autosuficiente.

\subsubsection{Transformación de Datos Tipados (TDT)}

La \textit{Transformación de Datos Tipados} (TDT)\parencite{jamalidinan2025tdt} representa
una línea más reciente de trabajo orientada a reorganizar datos de punto flotante antes
de aplicar compresión generalista. A diferencia de Bitshuffle, que opera sobre una
transposición fija de bits, TDT busca identificar posiciones de byte con comportamientos
estadísticos diferentes y reagruparlas para producir flujos más homogéneos.

La idea central es que, en una representación IEEE 754, no todos los bytes cumplen el
mismo papel ni presentan el mismo nivel de entropía. Los bytes asociados al exponente,
al signo o a distintas regiones de la mantisa pueden exhibir distribuciones distintas
según la naturaleza de los datos. TDT analiza esas diferencias a nivel de bloque y aplica
una reorganización reversible basada en características estadísticas, de modo que los
bytes con comportamiento semejante queden contiguos antes de ser entregados a un
compresor convencional.

Más que un códec independiente, TDT debe entenderse como un preprocesamiento
tipado orientado a mejorar la eficiencia de compresores generales sobre datos
numéricos. Su relevancia dentro del estado del arte radica en mostrar que la compresión
sin pérdida de flotantes no depende únicamente de predicción o de compresión
entrópica, sino también de la capacidad de reordenar la representación binaria para
hacer más visible su estructura estadística.

Esta línea de trabajo es especialmente pertinente para el problema de esta tesis porque
el vector de estado de un simulador Schr\"odinger se almacena como una secuencia de
amplitudes complejas en doble precisión, usualmente intercaladas como pares
real/imaginario. En ese formato, la compresibilidad no depende solo de la magnitud de
las amplitudes, sino también de patrones binarios en exponentes y mantisas. Por ello,
filtros reversibles como \emph{shuffle} o \emph{bitshuffle} resultan relevantes: pueden
agrupar bytes o bits homólogos de valores IEEE 754 y exponer regularidades que no son
visibles si cada amplitud se trata como un bloque opaco de 16 bytes.
